% GENERATED FILE. Edit canonical lessons, then run npm run generate:books.
% canonical-source-hash: fnv1a64:5b1317ffde503fb7
% canonical-lessons: ZH-C467-piqi, ZH-C467-gexing, ZH-W467-xing-nature, ZH-C467-tianxing, ZH-C467-kailang, ZH-C467-huopo

\chapter{What Kind of Person}
\label{ch:zh-what-kind-of-person}
\hlchaptermodality{467}

\begin{chapteropening}
\textbf{By the end of this chapter:} \emph{I can talk about someone's temper, personality and nature, and say someone is cheerful or lively.}

Complete the last lesson of chapter 467: I can talk about someone's temper, personality and nature, and say someone is cheerful or lively.
\end{chapteropening}

\section[píqi]{píqi — temper; temperament}

\label{lesson:ZH-C467-piqi}



\begin{quote}
Before the new one: say the Mandarin for to give a gentle smile, then the Mandarin for peaceful inside.
\end{quote}

\subsection*{You'll want to know: píqi}
\emph{píqi} — \textquotedblleft{}temper; temperament\textquotedblright{}.

\begin{quote}
\emph{wǒ bàba de píqi bǐ wǒ māma hǎo} — My dad is better-tempered than my mum.
\end{quote}

Its characters come later; for now it is a word for the ear.

\begin{cousinweb}
\emph{píqi} is built from \emph{pí} 'the spleen' and \emph{qì} 'breath', both inherited from Middle Chinese, \emph{bjie} and \emph{khj+jH} (Guangyun, in Baxter's transcription), which explains the first, after an older book, as the organ of earth that helps the stomach. Temper as the breath of the spleen is the usual explanation, not one the dictionary gives.
\end{cousinweb}

\begin{grammarlens}[title={What you've built}]
One more word for saying how it was.
\end{grammarlens}

\subsection*{Your turn}
\begin{itemize}
\raggedright
  \item \emph{Say it:} \emph{píqi}
  \item \emph{Say it:} \emph{píqi}, once more
  \item \emph{Say it:} \emph{wǒ bàba de píqi bǐ wǒ māma hǎo} — My dad is better-tempered than my mum.
  \item \emph{Recall:} say the Mandarin for to give a gentle smile, then the Mandarin for peaceful inside, then say \emph{píqi} again
\end{itemize}

\begin{tcolorbox}[breakable,colback=teal!4,colframe=teal!35!black,title={Before you move on}]
What does píqi mean? (\textquotedblleft{}Temper; temperament\textquotedblright{}.) Say it once more.
\end{tcolorbox}
\section[\texorpdfstring{\zh{个性} (gèxìng)}{个性 (gèxìng)}]{\zh{个性} (gèxìng) — individuality; a strong personality}

\label{lesson:ZH-C467-gexing}



\begin{quote}
Before the new one: say the Mandarin for peaceful inside, then the Mandarin for a person's temperament.
\end{quote}

\subsection*{You'll want to know: \zh{个性}}
\textbf{\zh{个性}} — \emph{gèxìng} — \textquotedblleft{}individuality; a strong personality\textquotedblright{}.

\begin{quote}
\emph{tā hěn yǒu gèxìng} — She has a lot of personality.
\end{quote}

The new character, in pieces: \zh{性} — eight strokes in two pieces: on the left the heart piece, a dot, a second dot and a long vertical; on the right \zh{生} \emph{shēng}, five strokes. The next lesson puts it in your hand.

\begin{cousinweb}
\textbf{\zh{个性}} (\emph{gèxìng}) is built from \zh{个} 'single' and \zh{性} 'nature', both inherited from Middle Chinese, \emph{kaH} and \emph{sjengH} (Guangyun, in Baxter's transcription), which glosses the first 'one-sided' and the second 'nature and conduct'.
\end{cousinweb}

\begin{grammarlens}[title={What you've built}]
One more word for saying how it was.
\end{grammarlens}

\subsection*{Your turn}
\begin{itemize}
\raggedright
  \item \emph{Say it:} \emph{gèxìng}
  \item \emph{Say it:} \emph{gèxìng}, once more
  \item \emph{Say it:} \emph{tā hěn yǒu gèxìng} — She has a lot of personality.
  \item \emph{Recall:} say the Mandarin for peaceful inside, then the Mandarin for a person's temperament, then say \emph{gèxìng} again
\end{itemize}

\begin{tcolorbox}[breakable,colback=teal!4,colframe=teal!35!black,title={Before you move on}]
What does \zh{个性} mean? (\textquotedblleft{}Individuality; a strong personality\textquotedblright{}.) Say it once more.
\end{tcolorbox}
\section[\texorpdfstring{\zh{性} (xìng)}{性 (xìng)}]{\zh{性} — the heart piece, then \zh{生}}

\label{lesson:ZH-W467-xing-nature}



\begin{quote}
The word first, then the character inside it.

\begin{itemize}
\raggedright
  \item \emph{Recall:} say \emph{gèxìng} — \textbf{R1}, one lesson back, before it decays
\end{itemize}

\begin{quote}
\zh{生} \zh{怕}
\end{quote}

\emph{Read it:} each one aloud, with nothing said first — \emph{shēng}, \emph{pà}
\end{quote}

\subsection*{Script — the heart piece, then \zh{生}}
\hlblockfigure{\detokenize{figures/ZH-W467-xing-nature-filmstrip.pdf}}{How \zh{性} is written, stroke by stroke}

\begin{quote}
\zh{性}
\end{quote}

\textbf{Two pieces, left before right.} Count the pieces, not the strokes.

\textbf{The left piece, three strokes: the heart piece.} A short dot on the left, a short dot on the right, then a long vertical down between them.

\textbf{The right piece, \zh{生}, five strokes.} A short stroke falling to the left at the top, a bar, a second bar, a vertical down through them, then a long bar across the foot.

3 + 5 = eight.

\subsection*{Your turn}
\begin{itemize}
\raggedright
  \item \emph{Write it:} \zh{性}, all eight, with the model in front of you
  \item \emph{Cover:} write it again from memory, and count the pieces, not the strokes
\end{itemize}

\begin{tcolorbox}[breakable,colback=teal!4,colframe=teal!35!black,title={Before you move on}]
How many pieces does \zh{性} have? (\textbf{Two: three strokes and five}.) Which written character is on the right? (\textbf{\zh{生}}.)

Source: \href{https://raw.githubusercontent.com/chanind/hanzi-writer-data/68d10a4b21150cae5e1ebbd223eed289cf32d90c/data/\%E6\%80\%A7.json}{Hanzi Writer Data \zh{性}, pinned stroke paths and medians}.
\end{tcolorbox}
\section[\texorpdfstring{\zh{天性} (tiānxìng)}{天性 (tiānxìng)}]{\zh{天性} (tiānxìng) — one's nature; an inborn trait}

\label{lesson:ZH-C467-tianxing}



\begin{quote}
Before the new one: say the Mandarin for a person's temperament, then the Mandarin for a strong personality.
\end{quote}

\subsection*{You'll want to know: \zh{天性}}
\textbf{\zh{天性}} — \emph{tiānxìng} — \textquotedblleft{}one's nature; an inborn trait\textquotedblright{}.

\begin{quote}
\textbf{\zh{安静是他的天性}} \emph{ānjìng shì tā de tiānxìng} — Being quiet is in his nature.
\end{quote}

\begin{cousinweb}
\textbf{\zh{天性}} (\emph{tiānxìng}) is built from \zh{天} 'heaven' and \zh{性} 'nature', both inherited from Middle Chinese, \emph{then} and \emph{sjengH} (Guangyun, in Baxter's transcription), which glosses the first 'the dark heaven above' and the second 'nature and conduct': the nature heaven gives.
\end{cousinweb}

\begin{grammarlens}[title={What you've built}]
One more word for saying how it was.
\end{grammarlens}

\subsection*{Your turn}
\begin{itemize}
\raggedright
  \item \emph{Say it:} \emph{tiānxìng}
  \item \emph{Say it:} \emph{tiānxìng}, once more
  \item \emph{Say it:} \emph{ānjìng shì tā de tiānxìng} — Being quiet is in his nature.
  \item \emph{Recall:} say the Mandarin for a person's temperament, then the Mandarin for a strong personality, then say \emph{tiānxìng} again
\end{itemize}

\begin{tcolorbox}[breakable,colback=teal!4,colframe=teal!35!black,title={Before you move on}]
What does \zh{天性} mean? (\textquotedblleft{}One's nature; an inborn trait\textquotedblright{}.) Say it once more.
\end{tcolorbox}
\section[kāilǎng]{kāilǎng — cheerful and outgoing}

\label{lesson:ZH-C467-kailang}



\begin{quote}
Before the new one: say the Mandarin for a strong personality, then the Mandarin for an inborn trait.
\end{quote}

\subsection*{You'll want to know: kāilǎng}
\emph{kāilǎng} — \textquotedblleft{}cheerful and outgoing\textquotedblright{}.

\begin{quote}
\emph{wǒ de péngyou lǐ tā zuì kāilǎng} — Of all my friends, she's the most cheerful.
\end{quote}

Its characters come later; for now it is a word for the ear.

\begin{cousinweb}
\emph{kāilǎng} is built from \zh{开} 'to open' and \emph{lǎng} 'bright', both inherited from Middle Chinese, \emph{khoj} and \emph{langX} (Guangyun, in Baxter's transcription), which glosses the second 'bright'. The sense of a cheerful nature is not given there.
\end{cousinweb}

\begin{grammarlens}[title={What you've built}]
One more word for saying how it was.
\end{grammarlens}

\subsection*{Your turn}
\begin{itemize}
\raggedright
  \item \emph{Say it:} \emph{kāilǎng}
  \item \emph{Say it:} \emph{kāilǎng}, once more
  \item \emph{Say it:} \emph{wǒ de péngyou lǐ tā zuì kāilǎng} — Of all my friends, she's the most cheerful.
  \item \emph{Recall:} say the Mandarin for a strong personality, then the Mandarin for an inborn trait, then say \emph{kāilǎng} again
\end{itemize}

\begin{tcolorbox}[breakable,colback=teal!4,colframe=teal!35!black,title={Before you move on}]
What does kāilǎng mean? (\textquotedblleft{}Cheerful and outgoing\textquotedblright{}.) Say it once more.
\end{tcolorbox}
\section[huópō]{huópō — full of life; vivacious}

\label{lesson:ZH-C467-huopo}



\begin{quote}
Before the new one: say the Mandarin for an inborn trait, then the Mandarin for cheerful and outgoing.
\end{quote}

\subsection*{You'll want to know: huópō}
\emph{huópō} — \textquotedblleft{}full of life; vivacious\textquotedblright{}.

\begin{quote}
\emph{wǒ nǚ'ér hěn huópō yě hěn ài xiào} — My daughter is very lively and laughs a lot.
\end{quote}

Its characters come later; for now it is a word for the ear.

\begin{cousinweb}
\emph{huópō} is built from \zh{活} 'alive' and \emph{pō} 'to splash'. The first is inherited from Middle Chinese \emph{hwat} (Guangyun, in Baxter's transcription), glossed 'not dead'; no dictionary we use records an older form of the second, so its history is not settled.
\end{cousinweb}

\begin{grammarlens}[title={What you've built}]
One more word for saying how it was.
\end{grammarlens}

\subsection*{Your turn}
\begin{itemize}
\raggedright
  \item \emph{Say it:} \emph{huópō}
  \item \emph{Say it:} \emph{huópō}, once more
  \item \emph{Say it:} \emph{wǒ nǚ'ér hěn huópō yě hěn ài xiào} — My daughter is very lively and laughs a lot.
  \item \emph{Recall:} say the Mandarin for an inborn trait, then the Mandarin for cheerful and outgoing, then say \emph{huópō} again
\end{itemize}

\begin{tcolorbox}[breakable,colback=teal!4,colframe=teal!35!black,title={Before you move on}]
What does huópō mean? (\textquotedblleft{}Full of life; vivacious\textquotedblright{}.) Say it once more.
\end{tcolorbox}
